\documentclass{article}
%  ╔══════════════════════════════════════════════════════════╗
%  ║                         Title                            ║
%  ╚══════════════════════════════════════════════════════════╝
%NOTE: I want \title,\author to be near top, and this \renewcommand{\maketitle} fails if these are above it.
\usepackage[utf8]{inputenc}         % Used to compile any UTF-8 Character (and supports Unicode)
\usepackage{titling}

\renewcommand{\maketitle}{
    \quad
    \vspace{1em}
  \begin{center}
      \LARGE\bfseries \thetitle \par
      \vspace{0.6em}
      \large
  \end{center}
    \vspace{1em}
}

\title{Some Harmonic Analysis}
\author{Nathanael Chwojko-Srawley}
\date{\today}

%  ╔══════════════════════════════════════════════════════════╗
%  ║                         packages                         ║
%  ╚══════════════════════════════════════════════════════════╝

\usepackage{extarrows}         % Used to compile any UTF-8 Character (and supports Unicode)
\usepackage{stackengine}
\usepackage{colortbl}
\usepackage{scalerel}
\usepackage{amsmath, amssymb} % Used to write better mathematical expressions (ex. \mathbb{R})
\usepackage{stackrel}         % for left-right arrows, staking symbol above and bellow
\usepackage{mathrsfs}
\usepackage{bbm}                     % for mathbbm{1}
\usepackage{euscript}
\usepackage{commutative-diagrams}
\usepackage[font=itshape]{quoting}
\usepackage{mathtools}
  \usepackage{enumitem}
\usepackage{amsthm}           % Used to create theorem enviroments.
\usepackage{graphicx}         % Let's you use images in LaTeX; ex the \includegraphics command.
\usepackage{hhline}
%\usepackage{luatexja}
%\usepackage{fontspec}
\usepackage{dsfont}
\usepackage{wasysym}          % emoji's!
\usepackage{float}            % package with ``H'' for figures and tables, ex. Images, tables, etc
%\usepackage{subcaption}      % More options with sub-captions.
\usepackage{setspace}         % More flexibility on spacing in document.
\usepackage{hyperref}         % let's me put hyperlinks
%\usepackage{tikz}             % for commutative diagrams https://tex.stackexchange.com/questions/419221/how-to-do-the-pushout-with-universal-property
\usepackage{tikz-cd}          % for commutative diagrams https://tex.stackexchange.com/questions/419221/how-to-do-the-pushout-with-universal-property
%\usepackage{quiver}           % for better commutative diagrams
\usepackage{bookmark}         % creating heading shortcut bar on the left as in word
\usepackage{longtable}        % create tables that extend past one page
%\usepackage{siunitx}         % required for alignment
%\usepackage{multirow}        % for multiple rows
\usepackage{microtype}        % makes nice text. Fixes some under-/over- filled box problems.
%\usepackage{fancyvrb}        % Let's you write code (different style than listings).
\usepackage[most]{tcolorbox} 	      % Makes nice boxes for thms, cor, prop, or anything else.
%\usepackage{enumitem}        % More flexibility on enumeration (don't know much).
\usepackage{xcolor}           % Colour text.
\usepackage{wrapfig}          % To wrap text around figure
\usepackage{listings}         % Create nice colour code blocks (customized in preamble)
\usepackage{thmtools}         % Customize theorem environment (in preamble).
%\usepackage{marvosym}        % Provide ``basic'' symbols (hazard, religious, smiley face etc.)
\usepackage{array}            % \begin{table}{>{\em}c c} -> 1st column italic (bfseries for bold)
%\usepackage[english]{babel}  % If I want to blindtext in english.
\usepackage{blindtext}
\usepackage[a4paper, total={6in, 8.5in}]{geometry}
%\usepackage[symbol]{footmisc} % For daggers and asterisk for footnotes. 1: *, 2: dagger, etc.
\usepackage{fancyhdr}
\usepackage{titlesec}
%\usepackage{pst-plot}          %To make nice plots: https://tex.stackexchange.com/questions/47594/how-can-i-make-a-graph-of-a-function#47596
% \usepackage[answerdelayed]{exercise}
\usepackage{etoolbox}          % for Dynkin diagrams
\usepackage{dynkin-diagrams}   % for Dynkin diagrams

% ╔═════════════════════════════════════════════════════╗
% ║           for figures via inkspace                  ║
% ╚═════════════════════════════════════════════════════╝

% to import figures via: https://github.com/gillescastel/inkscape-figures
\usepackage{import}
\usepackage{pdfpages}
\usepackage{transparent}


% This command is the ONLY command imported via the packages snippet, think of it as a tiny package written out
\newcommand{\incfig}[2][1]{%
    \def\svgwidth{#1\columnwidth}
    \import{./figures/}{#2.pdf_tex}
}

% ╔═══════════════════════════════════════════════════════════╗
% ║                 bibliography and citing                   ║
% ╚═══════════════════════════════════════════════════════════╝

\usepackage{csquotes}
% \usepackage[backend=biber,citestyle=alphabetic,bibstyle=authortitle]{biblatex}
\usepackage{biblatex}

\addbibresource{bibliography.bib}

% ╔═══════════════════════════════════════════════════════════╗
% ║                    colours (colors)                       ║
% ╚═══════════════════════════════════════════════════════════╝

\definecolor{dkgreen}{rgb}{0,0.6,0}
\definecolor{gray}{rgb}{0.5,0.5,0.5}
\definecolor{mauve}{rgb}{0.58,0,0.82}
\definecolor{lightyellow}{rgb}{1,1,0.6}
\definecolor{reddish}{HTML}{A01A3A}

%  ╔══════════════════════════════════════════════════════════╗
%  ║                       book format                        ║
%  ╚══════════════════════════════════════════════════════════╝

\pagestyle{fancy}

\setlength\headheight{20pt}

\renewcommand{\sectionmark}[1]{\markright{\thesection.\ #1}}

\fancyfoot[C,CO]{\textcolor{gray}{Nathanael Chwojko-Srawley}}
\fancyfoot[R,RO]{\thepage}{}

% Create a new page style for the first page
\fancypagestyle{firstpage}{
  \fancyhf{} % Clear header and footer
  \fancyhead[L]{\theauthor}
  \fancyhead[R]{\today} % Date in the top-right
  \fancyfoot[C]{\textcolor{gray}{Nathanael Chwojko-Srawley}}
  \fancyfoot[R]{\thepage}
}

% Apply the first page style conditionally
\AtBeginDocument{\thispagestyle{firstpage}}

%have format the title command
\newcommand{\chapter}[1]{%
  \clearpage
  \vspace*{30pt} % Space before the chapter title
  {\normalfont\huge\bfseries #1\par} % Chapter title formatting
  \vspace{20pt} % Space after the chapter title
}


\setlength{\parindent}{0pt} % don't like indentation infront of paragraphs
\setlength{\parskip}{1ex}   %so that there is still space between paragarphs


%  ╔══════════════════════════════════════════════════════════╗
%  ║                       Shortcuts                          ║
%  ╚══════════════════════════════════════════════════════════╝

%\ExerciseLevelInToc

% I like original mathcal better, but need lowercase.
\DeclareFontFamily{OT1}{pzc}{}
\DeclareFontShape{OT1}{pzc}{m}{it}{<-> s * [1.10] pzcmi7t}{}
\DeclareMathAlphabet{\mathpzc}{OT1}{pzc}{m}{it}


% mathbb
\newcommand{\A}{{\mathbb{A}}}
\newcommand{\C}{{\mathbb{C}}}
\newcommand{\F}{{\mathbb{F}}}
\newcommand{\N}{{\mathbb{N}}}
\newcommand{\Q}{{\mathbb{Q}}}
\newcommand{\R}{{\mathbb{R}}}
\newcommand{\V}{{\mathbb{V}}}
\newcommand{\Z}{{\mathbb{Z}}}
\newcommand{\BI}{{\mathbb{I}}}
\renewcommand{\H}{{\mathbb{H}}}
\renewcommand{\P}{{\mathbb{P}}}


% mathcal
% \newcommand{\co}{{\mathfrak{o}}} %mathcal{o} looks like wreath product, so I changed it to mathfrak
\newcommand{\co}{{\mathpzc{o}}}
\newcommand{\CA}{{\mathcal{A}}}
\newcommand{\CB}{{\mathcal{B}}}
\newcommand{\CC}{{\mathcal{C}}}
\newcommand{\CD}{{\mathcal{D}}}
\newcommand{\CF}{{\mathcal{F}}}
\newcommand{\CG}{{\mathcal{G}}}
\newcommand{\CH}{{\mathcal{H}}}
\newcommand{\CI}{{\mathcal{I}}}
\newcommand{\CK}{{\mathcal{K}}}
\newcommand{\CL}{{\mathcal{L}}}
\newcommand{\CM}{{\mathcal{M}}}
\newcommand{\CN}{{\mathcal{N}}}
\newcommand{\CO}{{\mathcal{O}}}
\newcommand{\CQ}{{\mathcal{Q}}}
\newcommand{\CR}{{\mathcal{R}}}
\newcommand{\CS}{{\mathcal{S}}}
\newcommand{\CT}{{\mathcal{T}}}
\newcommand{\CV}{{\mathcal{V}}}
\newcommand{\CZ}{{\mathcal{Z}}}
% EXCPETION:
\newcommand{\CP}{{\mathcal{P}}}
% \newcommand{\CP}{\mathbb{C}\mathrm{P}}
\newcommand{\RP}{\mathbb{R}\mathrm{P}}


%Euscript
\newcommand{\EB}{\EuScript{B}}
\newcommand{\EC}{{\EuScript{C}}}
\newcommand{\EF}{{\EuScript{F}}}
\newcommand{\EL}{{\EuScript{L}}}
\newcommand{\EO}{{\EuScript{O}}}


%mathfrak (lowercase)
\newcommand{\fa}{{\mathfrak{a}}}
\newcommand{\fb}{{\mathfrak{b}}}
\newcommand{\fc}{{\mathfrak{c}}}
\newcommand{\fd}{{\mathfrak{d}}}
\newcommand{\fm}{{\mathfrak{m}}}
\newcommand{\fn}{{\mathfrak{n}}}
\newcommand{\fo}{{\mathfrak{o}}}
\newcommand{\fp}{{\mathfrak{p}}}
\newcommand{\fq}{{\mathfrak{q}}}


%mathfrak (upper case)
\newcommand{\FA}{{\mathfrak{A}}}
\newcommand{\FB}{{\mathfrak{B}}}
\newcommand{\FC}{{\mathfrak{C}}}
\newcommand{\FD}{{\mathfrak{D}}}
\newcommand{\FK}{{\mathfrak{K}}}
\newcommand{\FM}{{\mathfrak{M}}}
\newcommand{\FN}{{\mathfrak{N}}}
\newcommand{\FO}{{\mathfrak{O}}}
\newcommand{\FP}{{\mathfrak{P}}}


%mathscr
\newcommand{\ff}{{\mathscr{F}}}
\newcommand{\SA}{\mathscr{A}}
\newcommand{\SC}{\mathscr{C}}
\newcommand{\SF}{\mathscr{F}}
\newcommand{\SG}{\mathscr{G}}
\newcommand{\SH}{\mathscr{H}}
\newcommand{\SI}{\mathscr{I}}
\newcommand{\SK}{\mathscr{K}}
\newcommand{\SN}{\mathscr{N}}
\newcommand{\SQ}{\mathscr{Q}}
\newcommand{\SR}{\mathscr{R}}
\newcommand{\SV}{\mathscr{V}}
\newcommand{\SZ}{\mathscr{Z}}


% operatorname -- 2 letters (lowercase and upper case)
\newcommand{\ab}{{\operatorname{ab}}}
\newcommand{\ad}{{\operatorname{ad}}}
\newcommand{\ch}{{\operatorname{ch}}}
\newcommand{\ev}{{\operatorname{ev}}}
\newcommand{\id}{{\operatorname{id}}}
\newcommand{\im}{{\operatorname{im}}}
\newcommand{\nm}{{\operatorname{Nm}}}
\newcommand{\op}{{\operatorname{op}}}
\newcommand{\rk}{{\operatorname{rk}}}
\newcommand{\sh}{{\operatorname{sh}}}
\newcommand{\tr}{{\operatorname{tr}}}

\newcommand{\Ab}{{\operatorname{Ab}}}
\newcommand{\Br}{{\operatorname{Br}}}
\newcommand{\Ch}{{\operatorname{Ch}}}
\newcommand{\Cl}{{\operatorname{Cl}}}
\newcommand{\GL}{{\operatorname{GL}}}
\newcommand{\Nm}{{\operatorname{Nm}}}
\newcommand{\SL}{{\operatorname{SL}}}

\renewcommand{\Re}{{\operatorname{Re}}}
\renewcommand{\Im}{{\operatorname{Im}}}


%categories
\newcommand{\Grp}{{\textbf{Grp}}}
\newcommand{\Fld}{{\textbf{Fld}}}
\newcommand{\Sh}{{\textbf{Sh}}}
\newcommand{\Set}{{\textbf{Set}}}
\newcommand{\Mod}{{\textbf{Mod}}}
\newcommand{\PSh}{{\textbf{PSh}}}
\newcommand{\Sch}{{\textbf{Sch}}}
\newcommand{\Top}{{\textbf{Top}}}
\newcommand{\RgSp}{{\textbf{RgSp}}}
\newcommand{\Ring}{{\textbf{Ring}}}


%operatorname -- 3+ letters (lowercase and upper case)
\newcommand{\rad}{{\operatorname{rad}}}
\newcommand{\sep}{{\operatorname{sep}}}
\newcommand{\res}{{\operatorname{res}}}
\newcommand{\sgn}{{\operatorname{sgn}}}
\newcommand{\pre}{{\operatorname{pre}}}
\newcommand{\ord}{{\operatorname{ord}}}
\newcommand{\vol}{{\operatorname{vol}}}
\newcommand{\lcm}{{\operatorname{lcm}}}

\newcommand{\conj}{{\operatorname{conj}}}
\newcommand{\disc}{{\operatorname{disc}}}
\newcommand{\stab}{{\operatorname{stab}}}
\newcommand{\kdim}{{\operatorname{kdim}}}
\newcommand{\diag}{{\operatorname{diag}}}

\newcommand{\trdeg}{\operatorname{trdeg}}
\newcommand{\coker}{{\operatorname{coker}}}
\newcommand{\codim}{{\operatorname{codim}}}

\newcommand{\Ann}{{\operatorname{Ann}}}
\newcommand{\Aut}{{\operatorname{Aut}}}
\newcommand{\Der}{{\operatorname{Der}}}
\newcommand{\Div}{{\operatorname{Div}}}
\newcommand{\Emb}{{\operatorname{Emb}}}
\newcommand{\End}{{\operatorname{End}}}
\newcommand{\Ext}{{\operatorname{Ext}}}
\newcommand{\Gal}{{\operatorname{Gal}}}
\newcommand{\Hom}{{\operatorname{Hom}}}
\newcommand{\Ind}{{\operatorname{Ind}}}
\newcommand{\Inn}{{\operatorname{Inn}}}
\newcommand{\Int}{{\operatorname{int}}}
\newcommand{\Irr}{{\operatorname{Irr}}}
\newcommand{\Jac}{{\operatorname{Jac}}}
\newcommand{\Mor}[1]{\operatorname{Mor}(\textbf{#1})}
\newcommand{\Nil}{{\operatorname{Nil}}}
\newcommand{\Obj}{{\operatorname{Obj}}}
\newcommand{\Out}{{\operatorname{Out}}}
\newcommand{\Pic}{{\operatorname{Pic}\,}}
\newcommand{\Rep}{{\operatorname{Rep}}}
\newcommand{\Res}{{\operatorname{Res}}}
\newcommand{\Spl}{{\operatorname{Spl}}}
\newcommand{\Syl}{{\operatorname{Syl}}}
\newcommand{\Sym}{{\operatorname{Sym}}}
\newcommand{\Tor}{{\operatorname{Tor}}}

\newcommand{\Char}{{\operatorname{char}}}
\newcommand{\Frac}{{\operatorname{Frac}}}
\newcommand{\Frob}{{\operatorname{Frob}}}
\newcommand{\Proj}{{\operatorname{Proj}\,}}
\newcommand{\RMod}{{}_R\operatorname{Mod}}
\newcommand{\SExt}{{\mathcal{E}\emph{xt}}}
\newcommand{\SHom}{{\emph{Hom}}}
\newcommand{\Span}{{\operatorname{span}}}
\newcommand{\Spec}{{\operatorname{Spec}\,}}
\newcommand{\Supp}{{\operatorname{Supp}}}
\newcommand{\Tens}[1]{#1\otimes --}
\newcommand{\fHom}[1]{\operatorname{Hom}(#1, --)}

\newcommand{\mSpec}{{\operatorname{mSpec}}}
\newcommand{\cofHom}[1]{\operatorname{Hom}(--, #1)}


% connectors
% \renewcommand{\mid}{\big|}
\newcommand{\sse}{\subseteq}
\newcommand{\subg}{\leqslant}
\newcommand{\supg}{\geqslant}
\newcommand{\ssne}{\subsetneq}
\newcommand{\isosubg}{\lesssim}
\newcommand{\nsse}{\not\subseteq}
\newcommand{\nsubg}{\trianglelefteq}
\newcommand{\nsupg}{\trianglerighteq}
\newcommand{\csubg}{\blacktriangleleft}
\renewcommand{\mid}{\big|}


% Arrows
\newcommand{\rw}{\rightarrow}
\newcommand{\Rw}{\Rightarrow}
\newcommand{\lw}{\leftarrow}
\newcommand{\Lw}{\Leftarrow}
\newcommand{\lrw}{\leftrightarrow}
\newcommand{\LRw}{\Leftrightarrow}
\newcommand{\acton}{\curvearrowright}
\newcommand{\actson}{\curvearrowright}
\newcommand\mapsfrom{\mathrel{\reflectbox{\ensuremath{\mapsto}}}}

% arrows for commutative diagram: create four new angled double-struck arrows
\newcommand\myrot[1]{\mathrel{\rotatebox[origin=c]{#1}{$\Rightarrow$}}}
\newcommand\NEarrow{\myrot{45}}
\newcommand\NWarrow{\myrot{135}}
\newcommand\SWarrow{\myrot{-135}}
\newcommand\SEarrow{\myrot{-45}}


%shorthands
\newcommand{\oln}[1]{{\overline{#1}}}
\renewcommand{\bf}[1]{\textbf{#1}}
\newcommand{\qn}{\quad\newline}


% limits
\newcommand{\Lim}{{\varprojlim}}
\newcommand{\coLim}{{\varinjlim}}
\newcommand{\Dlim}{\lim\limits_{\longrightarrow}}
\newcommand{\coDlim}{\lim\limits_{\longleftarrow}}


%for saturated ideals in the localization section (I don't like these, I think I'll delete)
\newcommand{\LIm}{{}^e }
\newcommand{\LPre}{{}^c }

%  ╔══════════════════════════════════════════════════════════╗
%  ║                    Custom Commands                       ║
%  ╚══════════════════════════════════════════════════════════╝

% swap varphi and phi
\let\temp\phi
\let\phi\varphi
\let\varphi\temp

\newcommand{\placeholder}{\_\_\_\_\_}

\newcommand{\set}[2]{\left\{#1 \ \middle|\ #2\right\}}

%mod should not have the crazy amount of space to its right!
\renewcommand{\mod}{\,\operatorname{mod}\,}

% dcups
\newcommand{\dcup}{\sqcup}
\newcommand{\Dcup}{\coprod}
% \newcommand{\Dcup}{\bigsqcup}

\newcommand{\nbot}{\mathrel{\rlap{$\bot$}\mkern2mu/}}

%a nice large circle
\newcommand{\tikzcircle}[2][red,fill=black]{\tikz[baseline=-0.5ex]\draw[#1,radius=#2] (0,0) circle ;}%

\makeatletter
\newcommand*\bigcdot{\mathpalette\bigcdot@{.5}}
\newcommand*\bigcdot@[2]{\mathbin{\vcenter{\hbox{\scalebox{#2}{$\m@th#1\bullet$}}}}}
\makeatother

%somehow not a default command
\def\rddots{\mathstrut^{.^{.^{.}}}}

%% new delimiter
\DeclarePairedDelimiter{\ceil}{\lceil}{\rceil}


%  ╔══════════════════════════════════════════════════════════╗
%  ║                     environments                         ║
%  ╚══════════════════════════════════════════════════════════╝


%remember the option: breakable

% create tcolor boxes
\newtcbtheorem[number within=section]{thm}{Theorem}%
{
colback=green!5,
colframe=green!35!black,
fonttitle=\bfseries,
label type=theorem
}{th}

\newtcbtheorem[number within=section, use counter from=thm]{lem}{Lemma}%
{
colback=blue!5,
colframe=black!35!black,
fonttitle=\bfseries,
label type=lemma
}{lm}
\newtcbtheorem[number within=section, use counter from=thm]{prop}{Proposition}%
{
colback=white!5,
colframe=white!35!black,
fonttitle=\bfseries,
label type=proposition
}{pr}
\newtcbtheorem[number within=section, use counter from=thm]{cor}{Corollary}%
{colback=blue!5,
colframe=blue!35!black,
fonttitle=\bfseries,
label type=corollary
}{co}
\newtcbtheorem[number within=section, use counter from=thm]{axiom}{Axiom}%
{colback=black!5,
colframe=yellow!35!black,
fonttitle=\bfseries,
label type=axiom
}{ax}
\newtcbtheorem[number within=section, use counter from=thm]{defn}{Definition}%
{colback=red!5,
colframe=red!35!black,
fonttitle=\bfseries,
label type=definition
}{df}


% for <s-cr> to create \item[$\LRw$] instead of \item
\newenvironment{equivEnumerate}
 {\begin{enumerate}
	}{
\end{enumerate}}



%Custom enviroments
 \newenvironment{note}
 {\begin{tcolorbox}[enhanced, sharp corners, colback=white , borderline={0.2pt}{0pt}{black}]
   \begin{quote}\textbf{Note}:
   }{
   \end{quote}\end{tcolorbox}}

 \newenvironment{intuition}
 {\begin{quote}\textbf{Intuition}:
   }{
   \end{quote}}

\newtcbtheorem[number within=section, use counter from=thm]{titledBox}{}{%
  enhanced,
  sharp corners,
  colback=white,
  colframe=black,
  borderline={0.2pt}{0pt}{black},
  left=2mm,
  right=2mm,
  top=2mm,
  bottom=2mm,
  title style={opacity=1},
  attach title to upper,
  before upper={\textbf{\tcbtitletext}\quad},
  coltitle=black,
  fonttitle=\bfseries,
  separator sign={\quad}
}{box}


%better example and proof environments
\def\exampletext{Example} % If English
\colorlet{colexam}{red!55!black} % Global example color


\newtcbtheorem[number within=section, use counter from=thm]{example}{Example}{% \newtcolorbox[use counter=testexample]{testexamplebox}{%
% Example Frame Start
empty,% Empty previously set parameters
title={\exampletext \thetcbcounter #1},% use \thetcbcounter to access the testexample counter text
% Attaching a box requires an overlay
attach boxed title to top left,
   % Ensures proper line breaking in longer titles
   minipage boxed title,
% (boxed title style requires an overlay)
boxed title style={empty,size=minimal,toprule=0pt,top=4pt,left=3mm,overlay={}},
coltitle=colexam,fonttitle=\bfseries,
before=\par\medskip\noindent,parbox=false,boxsep=0pt,left=3mm,right=0mm,top=2pt,breakable,pad at break=0mm,
   before upper=\csname @totalleftmargin\endcsname0pt, % Use instead of parbox=true. This ensures parskip is inherited by box.
% Handles box when it exists on one page only
overlay unbroken={\draw[colexam,line width=.5pt] ([xshift=-0pt]title.north west) -- ([xshift=-0pt]frame.south west); },
% Handles multipage box: first page
overlay first={\draw[colexam,line width=.5pt] ([xshift=-0pt]title.north west) -- ([xshift=-0pt]frame.south west); },
% Handles multipage box: middle page
overlay middle={\draw[colexam,line width=.5pt] ([xshift=-0pt]frame.north west) -- ([xshift=-0pt]frame.south west); },
% Handles multipage box: last page
overlay last={\draw[colexam,line width=.5pt] ([xshift=-0pt]frame.north west) -- ([xshift=-0pt]frame.south west); }%
}{ex}



\def\prooftext{Proof} % If English
\NewDocumentEnvironment{Proof}{ O{} }
{
    \colorlet{colexam}{black!55!black} % Global example color
    \newtcolorbox[number within=section]{testproofbox}{%
        empty,% Empty previously set parameters
        title={\emph{\prooftext} \thetcbcounter: #1},
        attach boxed title to top left,
        minipage boxed title,
        boxed title style={empty,size=minimal,toprule=0pt,top=4pt,left=3mm,overlay={}},
        coltitle=colexam,
        fonttitle=\bfseries,
        before=\par\medskip\noindent,
        parbox=false,
        boxsep=0pt,
        left=3mm,
        right=0mm,
        top=2pt,
        breakable,
        pad at break=0mm,
        before upper=\csname @totalleftmargin\endcsname0pt,
        % Removed height constraints
        % Added flexible width
        width=\dimexpr\textwidth-3mm\relax,
        % Modified overlays
        overlay unbroken={\draw[colexam,line width=.5pt] ([xshift=-0pt]title.north west) -- ([xshift=-0pt]frame.south west); },
        overlay first={\draw[colexam,line width=.5pt] ([xshift=-0pt]title.north west) -- ([xshift=-0pt]frame.south west); },
        overlay middle={\draw[colexam,line width=.5pt] ([xshift=-0pt]frame.north west) -- ([xshift=-0pt]frame.south west); },
        overlay last={\draw[colexam,line width=.5pt] ([xshift=-0pt]frame.north west) -- ([xshift=-0pt]frame.south west); },%
    }
    \begin{testproofbox}
}
{\end{testproofbox}\endlist}


%  ╔══════════════════════════════════════════════════════════╗
%  ║                    for Dynkin diagrams                   ║
%  ╚══════════════════════════════════════════════════════════╝

\def\row#1/#2!{#1_{\IfStrEq{#2}{}{n}{#2}} & \dynkin{#1}{#2}\\}
\newcommand{\tble}[1]{
   \renewcommand*\do[1]{\row##1!}
   \[
      \begin{array}{ll}\docsvlist{#1}\end{array}
   \]
}

%  ╔══════════════════════════════════════════════════════════╗
%  ║                         links                            ║
%  ╚══════════════════════════════════════════════════════════╝

\definecolor{LinkColour}{HTML}{A64A3B} % favourite
% \definecolor{LinkColour}{HTML}{8C3D32} % 2nd place
\hypersetup{
  colorlinks,
  citecolor=black,
  filecolor=magenta,
  linkcolor=LinkColour,
  urlcolor=blue,
}


%  ╔══════════════════════════════════════════════════════════╗
%  ║                          misc                            ║
%  ╚══════════════════════════════════════════════════════════╝

\stackMath %to be able to stack text in math mode
\makeindex

%import options: all, bibliography, book-formatting, booksSetting, customEnvironments, dynkin, hw-formatting, language-setup, newcommands, packages, preamble, proofs, settings, tcolorbox, test.svg,
\begin{document}
\maketitle

This is a direct extract from my algebra notes that I wanted to make a stand alone blog,
see~\cite{NateAlgebra}.

\section{*Application to Harmonic Analysis}

Let $G$ be a finite abelian group so that all complex irreducible representations are 1-dimensional, and hence
characters. Label
this collection $\hat{G} = \Hom_{\Ab}(G, \C^*)$. This is a group with point-wise multiplication.

\begin{prop}{Pontryagin Dual}{pontryaginDual}
	Let $G$ be a finite abelian group, $\hat{G} = \Hom_\Ab(G, \C^*)$ the group of characters under
	point-wise multiplication. Then there is an isomorphism:
	\[
		G \cong \widehat{G}
	\]
	Furthermore, there is a natural isomorphism:
	\[
		G \cong \hat{\hat{G}}
	\]
\end{prop}

If $G$ is infinite, it is rarely the case that $G \cong \widehat{G}$ (think $\widehat{L^p}\cong L^q$)

\begin{Proof}
	As $G$ if finite abelian,
	\[
		G \cong \Z/n_1\Z\times \Z/n_2\Z/ \times \cdots \times \Z/n_k\Z
	\]
	where $n_1 \mid n_2 \mid \cdots \mid n_k$. Choose generators, $a_1, ..., a_k$. Define the indicator
	characters on the generators:
	\[
		\chi_i(a_j) = \begin{cases}
			e^{\frac{2\pi i}{n_i}} & i = j\\
				1 & i \ne j
		\end{cases}
	\]
	These certainly generate $\hat{G}$. Define the map on the generators:
	\[
		\Phi(a_i) = \chi_i
	\]
	Clearly this is an isomorphism.

	For the double-dual, take the natural map:
	\[
		\Phi(g_i) = \ev_{g_i}
	\]
	where as usual $\ev_{g_i}(\chi) = \chi(g_i)$. This is easy to check that is an isomorphism, completing the
	proof.
\end{Proof}


Let $G$ have the discrete topology so that the set of continuous functions from $G$ to $\C$ is all set
functions. We want a $\C$-algebra isomorphism between $\C[G]$ and $\Hom_\Top(G, \C)$. These two sets are
clearly bijective via the natural mapping:
\[
	\sum_i a_ig_i \mapsto (g_i \mapsto a_i)
\]
And this map preserves addition if addition is defined point-wise in the hom-set:
\begin{align*}
	&\ \Phi\left(\sum_i a_ig_i + \sum_i b_ig_i\right)\\
	&=\Phi\left(\sum_i (a_i+b_i)g_i\right)\\
	  &\mapsto (g_i \mapsto a_i+b_i)\\
	  &= (g_i \mapsto a_i) + (g_i \mapsto b_i)\\
	  &= \Phi\left(\sum_i a_ig_i\right) + \Phi\left(\sum_i b_ig_i\right)
\end{align*}
But it does \emph{not} preserve multiplication if the operation is pointwise multiplication:
\begin{align*}
	&\ \Phi\left(\sum_{g \in G} a_ig\sum_{g\in G} b_ig\right)\\
	&=\Phi\left(\sum_{g\in G}\sum_{g_kg_\ell=g} (a_kb_\ell)g\right)\\
	  &\mapsto (g_i \mapsto \sum_{g_kg_\ell=g_i} (a_kb_\ell))\\
	  &\ne (g_i \mapsto a_ib_i)\\
	  &\ne (g_i \mapsto a_i) \cdot (g_i \mapsto b_i)\\
	  &= \Phi\left(\sum_i a_ig_i\right) \cdot \Phi\left(\sum_i b_ig_i\right)
\end{align*}
This is because the points $g_i$ act on each other. There is a natural multiplication function defined on
$\Hom_\Top(G, \C)$ called \emph{convolution}:

\index{Convolution}
\begin{defn}{Convolution}{convolution}
	Let $p,q \in \Hom_\Top(G, \C)$. Then their \emph{convolution} is define pointwise as:
	\[
		(p*q)(g) = \sum_{g_ig_j = g} p(g_i)q(g_j) = \sum_{h \in G}p(h)q(h^{-1}g)
	\]
\end{defn}

\begin{lem}{Isomorphism With Convolution}{isoWithConvolution}
	Let $G$ be a finite abelian group. Then:
	\[
    (\C[G], +, \cdot) \cong_{k\text{-}\bf{Alg}} (\Hom_\Top(G, \C), +, *)
	\]
	where the left hand side is the usual group ring and the right hand side has convolution as multiplication
\end{lem}

\begin{Proof}
	refer to the above calculations.
\end{Proof}

Now, $(\Hom_\Top(G, \C), +, *)$ is a $|G|$-dimensional $\C$ vector-space with the natural basis
given by the indicator function:
\[
	\delta_h(g) = \begin{cases}
		1 & g=h\\
		0 & g \ne h
	\end{cases}
\]
Namely for any $f: G \to \C$ there is a unique decomposition:
\[
	f = \sum_{h \in G} f(h)\delta_h \qquad \text{so that} \qquad f(g) = \sum_{h \in G} f(h)\delta_h(g)
	= f(g)\cdot 1 = f(g)
\]
The key feature of this basis is that it acts like evaluation when taking the inner product. Define:
\[
	\langle f,g \rangle = \sum_{h \in G} f(h)\overline{g(h)}
\]
Then $\{\delta_g\}_{g \in G}$ and:
\[
	\langle f, \delta_{g_i} \rangle = f(g_i) \cdot \overline{\delta_{g_i}(g_i)} = f(g_i) \cdot 1 = f(g_i)
	= a_i
\]
This inner product can be thought as expressly chosen so that the collection $\delta_{g_i}$ forms an
orthonormal basis, or conversely given this standard inner product, the indicator function will always be an
orthonormal basis. By Pontryagin duality, $G\cong \hat{G}$, but this group isomorphism does not
\emph{a-priori} preserve the
orthogonality of the basis given by $G$. The key insight of Fourier analysis is that it does!
\begin{prop}{Character Form Orthogonal Basis: Fourier Basis}{characterFormOrthogonalBasis}
	The characters $\widehat{G}$ form an orthogonal basis for $\Hom_{\Top}(G, \C)$. Normalizing by
	$1/\sqrt{|G|}$ forms an orthonormal basis.
\end{prop}

\begin{Proof}
	If $i = j$:
	\begin{align*}
		\langle \chi_i, \chi_i \rangle &= \sum_{g\in G} \chi_i(g)\overline{\chi_i(g)}\\
									   &= \sum_{g \in G} |\chi_i(g)|^2\\
									   &= \sum_{g \in G} 1^2\\
									   &= |G|
	\end{align*}
	if $i \ne j$, Then as $\overline{\chi_j} = \chi_j^{-1}$ (as we are mapping to the unit circle), we get
	that $\chi' = \chi_i\chi_j^{-1}$ is a non-trivial character, so there is some $h$ such that $\chi'(h) \ne
	1$. Let
	\[
		\langle \chi_i, \chi_j \rangle = S = \sum_{g \in G} \chi'(g)
	\]
	Then:
	\[
		\chi'(h) S = \chi'(h) \sum_{g \in G} \chi'(g) = \sum_{g \in G} \chi'(h)\chi'(g) = \sum_{g \in G} \chi'(h+g)
	\]
	As $h+G = G$, this is just a re-arranging of terms, hence:
	\[
		\chi'(h)S =\sum_{g \in G} \chi'(g) = S \qquad \text{so} \qquad S(\chi'(h)-1) = 0
	\]
	Now, by assumption $\chi'(h)\ne 1$, so $\chi'(h)-1\ne 0$, so as $\C$ is an integral domain it must be that
	$S = 0$, implying:
	\[
		\langle \chi_i, \chi_j \rangle = 0
	\]
	as we sought to show.
\end{Proof}

This inner product can be thought of as extracting the ``phase'' information.
\begin{example}{Phase-Space Basis}{phaseSpBasis}
Let $G = \Z/2\Z \times \Z/2\Z = \{(0,0), (1,0), (0,1), (1,1)\}$. A character $\chi: G \to \C^*$ is determined by its values on the generators $g_1=(1,0)$ and $g_2=(0,1)$. Since both generators have order 2, we must have $\chi(g_1)^2 = \chi(g_1^2) = \chi(0,0) = 1$ and likewise for $g_2$. Thus, the character's values on the generators must be either $1$ or $-1$. This gives $2\times 2=4$ possible characters, which we can index by pairs $(k_1, k_2) \in \{0,1\}^2$.

	The four characters $\chi_{k_1, k_2}$ are defined by their values on the generators:
	\[
		\chi_{k_1, k_2}(1,0) = (-1)^{k_1} \qquad \text{and} \qquad \chi_{k_1, k_2}(0,1) = (-1)^{k_2}
	\]
	From this, the value for any element $(a,b) \in G$ is $\chi_{k_1,k_2}(a,b) = (-1)^{k_1a + k_2b}$.

This gives the following character table, which explicitly lists the functions that form our new ``basis'':

\begin{center}
\begin{tabular}{c|rrrr}
	\textbf{$\hat{G}$} & \textbf{(0,0)} & \textbf{(1,0)} & \textbf{(0,1)} & \textbf{(1,1)} \\
	\hline
	$\chi_{00}$ & 1 & 1 & 1 & 1 \\
	$\chi_{10}$ & 1 & -1 & 1 & -1 \\
	$\chi_{01}$ & 1 & 1 & -1 & -1 \\
	$\chi_{11}$ & 1 & -1 & -1 & 1 \\
\end{tabular}
\end{center}

We can verify the orthogonality relation from proposition~\ref{pr:characterFormOrthogonalBasis} directly. Let us
compute $\langle \chi_{10}, \chi_{11} \rangle$:
\begin{align*}
	\langle \chi_{10}, \chi_{11} \rangle &= \sum_{g \in G} \chi_{10}(g)\overline{\chi_{11}(g)} \\
	&= \chi_{10}(0,0)\chi_{11}(0,0) + \chi_{10}(1,0)\chi_{11}(1,0) + \chi_{10}(0,1)\chi_{11}(0,1) + \chi_{10}(1,1)\chi_{11}(1,1) \\
	&= (1)(1) + (-1)(-1) + (1)(-1) + (-1)(1) \\
	&= 1 + 1 - 1 - 1 = 0
\end{align*}

\end{example}

\index{Fourier Transform}\index{Gelfand Transform}
\begin{defn}{Fourier Transform (Gelfand Transform)}{fourierTransform(gelfandTransform)}
	Let $G$ be a finite abelian group with the discrete topology. Then the \emph{Fourier transform} is:
	\[
		(\widehat{-}): \Hom_\Top(G, \C) \mapsto \Hom_\Top(\hat{G}, \C) \qquad f \mapsto (\chi \mapsto \langle
		f, \chi\rangle)
	\]
	where the image of $f$ is often written as:
	\[
		\hat{f}(\chi) = \langle f, \chi \rangle = \sum_{g\in G} f(g)\overline{\chi(g)}
	\]
\end{defn}

The name \emph{Gelfand transform} comes from the fact that the Fourier transform is traditionally about real
periodic functions, or Legesgue measurable functions. It wasn't until later where mathematicians like Folland
and Gelfand generalized it to the context of locally compact abelian groups, where the above is the special
case where $G$ is finite.

\index{Fourier Reciprocity}
\begin{thm}{Fourier Reciprocity}{fourierReciprocity}
	The map given in definition~\ref{df:fourierTransform(gelfandTransform)} is a unitary $\C$-algebra isomorphism, with
	point-wise multiplication on $(\Hom_\Top(\hat{G}, \C), +, \cdot)$.
\end{thm}

Recall that unitary means $\langle \hat{f}, \hat{g} \rangle = \langle f,g \rangle$.
Often, the $\C$-algebra homomorphism is called the  \emph{Convolution Theorem}, the unitary property is called the
\emph{Plancherel's Theorem}, and invertibility (i.e. bijectivity) is called the \emph{Fourier Inversion
Theorem}.
\begin{Proof}
	\begin{enumerate}
		\item \textit{1. Linearity:} Let $f, h \in \Hom_\Top(G, \C)$ and $a,b \in \C$.
		\begin{align*}
			\widehat{af+bh}(\chi) &= \langle af+bh, \chi \rangle && \text{by definition} \\
			&= \sum_{g\in G} (af+bh)(g) \overline{\chi(g)} \\
			&= a\sum_{g\in G} f(g)\overline{\chi(g)} + b\sum_{g\in G} h(g)\overline{\chi(g)} && \text{by linearity of sums}\\
			&= a\langle f, \chi \rangle + b\langle h, \chi \rangle \\
			&= a\hat{f}(\chi) + b\hat{h}(\chi)
		\end{align*}
		Since this holds for all $\chi \in \hat{G}$, the transform is a linear map.

		\textit{2. $\C$-algebra Homomorphism (The Convolution Theorem):} We must show that the Fourier transform converts convolution into pointwise multiplication.
		\begin{align*}
			\widehat{f*h}(\chi) &= \langle f*h, \chi \rangle \\
			&= \sum_{k \in G} (f*h)(k) \overline{\chi(k)} && \text{by definition of inner product} \\
			&= \sum_{k \in G} \left( \sum_{g \in G} f(g)h(g^{-1}k) \right) \overline{\chi(k)} && \text{by definition of convolution} \\
			&= \sum_{g \in G} f(g) \sum_{k \in G} h(g^{-1}k) \overline{\chi(k)} && \text{swapping order of summation}
		\end{align*}
		Let $\ell  = g^{-1}k$, so $k = g\ell $. As $k$ runs over $G$, so does $\ell $. The inner sum becomes:
		\[
			\sum_{\ell  \in G} h(\ell ) \overline{\chi(g\ell )} = \sum_{\ell  \in G} h(\ell ) \overline{\chi(g)}\overline{\chi(\ell )} = \overline{\chi(g)} \sum_{\ell  \in G} h(\ell )\overline{\chi(\ell )} = \overline{\chi(g)} \hat{h}(\chi)
		\]
		Substituting this back into the main expression:
		\begin{align*}
			\widehat{f*h}(\chi) &= \sum_{g \in G} f(g) \left(\overline{\chi(g)} \hat{h}(\chi)\right) \\
			&= \hat{h}(\chi) \sum_{g \in G} f(g) \overline{\chi(g)} \\
			&= \hat{h}(\chi) \hat{f}(\chi)
		\end{align*}
		Thus, $\widehat{f*h} = \hat{f} \cdot \hat{h}$, proving the map is an algebra homomorphism.

		\textit{3. Isomorphism (The Fourier Inversion Theorem):} We show the map is invertible by constructing its inverse. Let $\mathcal{F}(f) = \hat{f}$. We claim the inverse is given by $\mathcal{F}^{-1}(\hat{f})(g) = \frac{1}{|G|}\sum_{\chi \in \hat{G}} \hat{f}(\chi)\chi(g)$.
		\begin{align*}
			\frac{1}{|G|}\sum_{\chi \in \hat{G}} \hat{f}(\chi)\chi(g) &= \frac{1}{|G|}\sum_{\chi \in \hat{G}} \left( \sum_{h \in G} f(h)\overline{\chi(h)} \right) \chi(g) && \text{by definition of } \hat{f}\\
			&= \frac{1}{|G|} \sum_{h \in G} f(h) \sum_{\chi \in \hat{G}} \chi(g)\overline{\chi(h)} && \text{swapping order of summation} \\
			&= \frac{1}{|G|} \sum_{h \in G} f(h) \sum_{\chi \in \hat{G}} \chi(g-h)
		\end{align*}
		The inner sum is an example of the \emph{Second Orthogonality Relation} and is a standard exercise the
		reader should do. In particular, the sum of all characters evaluated at a group element (giving all
		powers of a certain root of unity) is $|G|$ if $x=e$ and $0$ otherwise. Thus $\sum_{\chi \in \hat{G}} \chi(g-h) = |G|\delta_{g,h}$.
		\[
			\frac{1}{|G|} \sum_{h \in G} f(h) (|G|\delta_{g,h}) = \sum_{h \in G} f(h)\delta_{g,h} = f(g)
		\]
		Since we can recover $f$ uniquely from $\hat{f}$, the map is an isomorphism.

		\textit{4. Unitary Property (Plancherel's Theorem):} Note that here is where it is important for the
		inner product to be normalized. Let us calculate the inner product for the given definition.
		\begin{align*}
			\langle \hat{f}, \hat{h} \rangle_{\hat{G}} &= \sum_{\chi \in \hat{G}} \hat{f}(\chi) \overline{\hat{h}(\chi)} \\
			&= \sum_{\chi \in \hat{G}} \left(\sum_{g \in G} f(g)\overline{\chi(g)}\right) \overline{\left(\sum_{k \in G} h(k)\overline{\chi(k)}\right)} \\
			&= \sum_{\chi \in \hat{G}} \left(\sum_{g \in G} f(g)\overline{\chi(g)}\right) \left(\sum_{k \in G} \overline{h(k)}\chi(k)\right) \\
			&= \sum_{g \in G}\sum_{k \in G} f(g)\overline{h(k)} \sum_{\chi \in \hat{G}} \chi(k)\overline{\chi(g)} && \text{re-arranging terms}\\
			&= \sum_{g \in G}\sum_{k \in G} f(g)\overline{h(k)} (|G|\delta_{g,k}) && \text{by Second Orthogonality} \\
			&= |G|\sum_{g \in G} f(g)\overline{h(g)} \\
			&= |G|\langle f, h \rangle_G
		\end{align*}
		This proves \emph{Plancherel's Theorem}. To make the map strictly unitary, one must normalize the
		transform. If we define the unitary transform $\mathcal{U}(f) = \frac{1}{\sqrt{|G|}}\hat{f}$, then it
		follows immediately that $\langle \mathcal{U}(f), \mathcal{U}(h) \rangle = \langle f,h \rangle$,
		completing the proof.
	\end{enumerate}
\end{Proof}


\begin{example}{Fourier Transform}{fourierTransform}
	Take $G= \Z/2\Z\times \Z/2\Z$ form example~\ref{ex:phaseSpBasis}. Let us now compute the Fourier transform of a concrete function. Let $f = \delta_{(1,0)}$, so that $f(1,0) = 1$ and is zero elsewhere. This is a basis vector in the standard ``position'' basis. We find its coordinates in the character basis by computing $\hat{f}(\chi) = \langle f, \chi \rangle$ for each character.
\begin{align*}
	\hat{f}(\chi) = \langle \delta_{(1,0)}, \chi \rangle &= \sum_{g \in G} \delta_{(1,0)}(g)\overline{\chi(g)} \\
	&= 1 \cdot \overline{\chi(1,0)} + 0 + 0 + 0 = \chi(1,0)
\end{align*}
Using the character table, we can compute this value for each $\chi$:
\begin{itemize}
	\item $\hat{f}(\chi_{00}) = \chi_{00}(1,0) = 1$
	\item $\hat{f}(\chi_{10}) = \chi_{10}(1,0) = -1$
	\item $\hat{f}(\chi_{01}) = \chi_{01}(1,0) = 1$
	\item $\hat{f}(\chi_{11}) = \chi_{11}(1,0) = -1$
\end{itemize}
Thus, the original function $f = \delta_{(1,0)}$, which is a spike at one point in the ``position'' basis, has been transformed into a new function $\hat{f}$ on $\hat{G}$ with values $(1, -1, 1, -1)$.

Recall the inversion formula from the proof of theorem~\ref{th:fourierReciprocity}: $f(g) = \frac{1}{|G|}\sum_{\chi
\in \hat{G}}\hat{f}(\chi)\chi(g)$ (here $|G|=4$). Let's check the value at $g=(1,0)$:
\begin{align*}
	f(1,0) &= \frac{1}{4}\left( \hat{f}(\chi_{00})\chi_{00}(1,0) + \hat{f}(\chi_{10})\chi_{10}(1,0) + \hat{f}(\chi_{01})\chi_{01}(1,0) + \hat{f}(\chi_{11})\chi_{11}(1,0) \right) \\
	&= \frac{1}{4}\left( (1)(1) + (-1)(-1) + (1)(1) + (-1)(-1) \right) \\
	&= \frac{1}{4}(1 + 1 + 1 + 1) = 1
\end{align*}
Now let's check the value at $g=(0,1)$, where we expect 0:
\begin{align*}
	f(0,1) &= \frac{1}{4}\left( \hat{f}(\chi_{00})\chi_{00}(0,1) + \hat{f}(\chi_{10})\chi_{10}(0,1) + \hat{f}(\chi_{01})\chi_{01}(0,1) + \hat{f}(\chi_{11})\chi_{11}(0,1) \right) \\
	&= \frac{1}{4}\left( (1)(1) + (-1)(1) + (1)(-1) + (-1)(-1) \right) \\
	&= \frac{1}{4}(1 - 1 - 1 + 1) = 0
\end{align*}
\end{example}

Let us now connect this to Representation Theory. So far, we have established that for a finite abelian group
$G$, the Fourier transform is a unitary $\C$-algebra isomorphism between the algebra of functions on $G$ with
convolution, $(\Hom_\Top(G, \C), *)$, and the algebra of functions on the dual group $\hat{G}$ with pointwise
multiplication. The representation theory perspective shows that Fourier
analysis is the simplest case of a Artin-Wedderburn decomposition.
The key is to view the $\C$-group algebra $\C[G]$ with the natural $G$-action given by left multiplication
(convolution), namely $\rho: G \to \GL(\C[G])$ given by:
\[
	\rho(g): f \mapsto g * f
\]
From this viewpoint, the study of the $\C$-algebra $(\C[G], +, *)$ and the $\C[G]$-module $\C[G]$  is the same
thing that is the study of the regular
representation of $G$. Now, this representation can be broken down into the direct sum of 1-dimensional
irreducible representations. Namely, for a finite abelian group, the set of irreducible representations is
exactly the dual group: $\Irr(G) = \widehat{G}$ and so we get:
\[
	\C[G] \cong \bigoplus_{\chi \in \widehat{G}} V_\chi
\]
where each $V_\chi$ is a 1-dimensional subspace corresponding to the character $\chi$.

This decomposition is made concrete by the Fourier transform. The Fourier Inversion Formula, $f =
\frac{1}{|G|} \sum_{\chi \in \hat{G}} \hat{f}(\chi)\chi$, is the literal expression of this decomposition. It
writes the function $f$ as a sum of components, where each term $\hat{f}(\chi)\chi$ is the projection of $f$
onto the 1-dimensional irreducible subspace $V_\chi$. Then the Fourier coefficient $\hat{f}(\chi)$ is the
coordinate of $f$ in the irreducible component corresponding to $\chi$. Note that we can define projection functions $p_\chi = \frac{1}{|G|}\sum_{g \in G}\overline{\chi(g)}g$ that we
can construct in $\C[G]$ that project onto the irreducible subspace $V_\chi$. Convolving any function $f$ with
$p_\chi$ annihilates all other ``frequency components'', leaving the part of $f$ that lies in $V_\chi$. Hence, the Fourier basis \emph{the} unique basis that diagonalizes the natural action of the group on its own algebra of functions.

If $G$ is a \emph{non-abelian} finite group, the concept of a dual group $\hat{G}$ and a simple transform
onto it ceases to be useful (it is now a set with no natural group structure). However, the representation
theory perspective remains: the generalized Fourier transform (given by the \emph{Peter-Weyl Theorem})
decomposes a function $f$ on the group into a sum of \emph{matrix coefficients} corresponding to these
higher-dimensional matrix representations.


\subsection{Generalization to Locally Compact Abelian Groups}
\label{sec:lcaGeneralization}

What happens when we wish to analyze functions on abelian groups like $(\R, +)$, $(\Z, +)$, or the circle group $(S^1, \cdot)$? These are examples of \emph{Locally Compact Abelian (LCA)} groups. Moving from the finite to the infinite setting requires us to replace our finite sums with integrals, but to integrate on an arbitrary group, we first need a notion of ``volume'' or ``measure''.
The essential new ingredient is the \emph{Haar measure}. For any LCA group $G$, there exists a measure $\mu$ that gives a consistent way to integrate functions on the group.

\index{Haar Measure}
\begin{defn}{Haar Measure}{haarMeasure}
	A \emph{Haar measure} on an LCA group $G$ is a measure $\mu$ on the Borel $\sigma$-algebra of $G$ that is left-translation-invariant. That is, for any measurable set $E \subseteq G$ and any element $g \in G$:
	\[
		\mu(gE) = \mu(E)
	\]
	This measure is unique up to a positive scaling constant. For abelian groups, it is also right-translation-invariant.
\end{defn}
See~\cite{NateRealAnalysis} for details on how to construct such a measure on any LCA that is unique up to
scaling. This measure is the direct generalization of taking the order of the group.
\begin{example}{Haar Measure}{haarMeas}
\begin{enumerate}
	\item For a finite group with the discrete topology, the Haar measure is the counting measure, and $\mu(G) = |G|$.
	\item For $(\R, +)$, the Haar measure is the standard Lebesgue measure $dx$.
	\item For $(S^1, \cdot)$, the Haar measure can be taken as $d\theta$.
	\item For $(\Z, +)$, the Haar measure is again the counting measure.
\end{enumerate}
\end{example}

With the Haar measure, we can redefine our main objects of study by replacing every sum over $G$ with an integral with respect to $d\mu$.

\begin{enumerate}
	\item \emph{The Group Algebra $L^1(G)$:} The space of all functions becomes too large. The natural replacement for the convolution algebra $\C[G]$ is the space of absolutely integrable functions, $L^1(G, d\mu)$. For two functions $f, h \in L^1(G)$, their convolution is defined by the integral:
	\[
		(f*h)(g) = \int_G f(x) h(x^{-1}g) d\mu(x)
	\]
	This space $(L^1(G), +, *)$ forms a Banach algebra.

	\item \emph{The Dual Group $\widehat{G}$:} It was proven that the correct codomain should be $S^1$ (see
		\textcolor{red}{ref:HERE} for the details, Poincare duality paper). The dual group is the set of all \emph{continuous} group homomorphisms into $S^1$.
	\[
		\widehat{G} = \Hom_{\Top}(G, S^1)
	\]
	The dual group $\widehat{G}$ is itself an LCA group. It is crucial to note that, unlike the finite case, it is rarely true that $G \cong \widehat{G}$. For example:
	\[
		\widehat{\R} \cong \R, \qquad \widehat{\Z} \cong S^1, \qquad \widehat{S^1} \cong \Z
	\]
	This highlights why it is essential to treat $\hat{G}$ as a distinct space. The theory that governs this relationship is \emph{Pontryagin Duality}, which states that the double-dual is canonically isomorphic to the original group: $\widehat{\widehat{G}} \cong G$.

	\item \emph{The Fourier Transform:} The definition of the Fourier transform is the direct analogue of the finite case, replacing the sum with an integral. For a function $f \in L^1(G)\cap L^2(G)$, its Fourier transform is a function $\hat{f}$ on the dual group $\hat{G}$:
	\[
		\hat{f}(\chi) = \int_G f(g) \overline{\chi(g)} d\mu(g)
	\]
	Similarly, the inner product that defines the Hilbert space structure on $L^2(G)$ is:
	\[
		\langle f, h \rangle = \int_G f(g) \overline{h(g)} d\mu(g)
	\]
\end{enumerate}

With these definitions, the entire suite of theorems proven in the finite case carries over:

\begin{thm}{Fourier Reciprocity on LCA Groups}{fourierReciprocityLCA}
	Let $G$ be an LCA group with Haar measure $\mu$. Then there exists a unique Haar measure $\nu$ on the dual group $\widehat{G}$ such that the Fourier transform, appropriately normalized, is a unitary $\C$-algebra isomorphism. The main results are:
	\begin{enumerate}
		\item \emph{The Fourier Inversion Theorem:} For a sufficiently well-behaved function $f$ (e.g., $f \in L^1 \cap L^2$ such that $\hat{f} \in L^1$), the original function can be recovered by an integral over the dual group:
		\[
			f(g) = \int_{\widehat{G}} \hat{f}(\chi) \chi(g) d\nu(\chi)
		\]
		\item \emph{Plancherel's Theorem:} The Fourier transform extends uniquely from $L^1 \cap L^2$ to a unitary isomorphism of Hilbert spaces, $\mathcal{F}: L^2(G, \mu) \to L^2(\widehat{G}, \nu)$. This means it preserves the inner product:
		\[
			\langle \hat{f}, \hat{h} \rangle_{L^2(\widehat{G})} = \langle f, h \rangle_{L^2(G)}
		\]
		\item \emph{The Convolution Theorem:} For $f, h \in L^1(G)$, the Fourier transform maps convolution to pointwise multiplication:
		\[
			\widehat{f * h}(\chi) = \hat{f}(\chi) \cdot \hat{h}(\chi)
		\]
	\end{enumerate}
\end{thm}

\begin{Proof}
	see \textcolor{red}{cite:HERE} (it's not in my analysis book, and I didn't type it up in Pontryagin notes)
\end{Proof}

The correspondence between the finite and general LCA settings can be summarized as follows:

\begin{center}
\begin{tabular}{l|l}
	\textbf{Finite Abelian Case} & \textbf{LCA Case} \\
	\hline
	Group $G$ & LCA Group $G$ \\
	Order $|G|$ & Haar Measure $\mu$ \\
	Group Algebra $\C[G]$ & Banach Algebra $L^1(G, \mu)$ \\
	Hilbert Space $\Hom(G, \C)$ & Hilbert Space $L^2(G, \mu)$ \\
	Summation $\sum_{g \in G}$ & Integration $\int_G d\mu(g)$ \\
	Dual Group $\widehat{G} = \Hom(G, \C^\times)$ & Dual Group $\widehat{G} = \Hom_\Top(G, S^1)$ \\
	Inner Product $\sum f(g)\overline{h(g)}$ & Inner Product $\int f(g)\overline{h(g)}d\mu(g)$ \\
	Fourier Transform $\sum f(g)\overline{\chi(g)}$ & Fourier Transform $\int f(g)\overline{\chi(g)}d\mu(g)$
\end{tabular}
\end{center}

Further interesting facts can be found here:


\url{https://mathoverflow.net/questions/37021/is-fourier-analysis-a-special-case-of-representation-theory-or-an-analogue}

\subsection{The Algebraic-Geometric Meaning of the Dual Group}

In algebraic geometry, a fundamental principle is that the geometry of a space can be recovered from its algebra of functions. For an affine space $k^n$, the corresponding algebra is the polynomial ring $k[x_1, \dots, x_n]$. The points of the space correspond to the maximal ideals of this algebra. Specifically, for any point $p=(a_1, \dots, a_n) \in k^n$, the evaluation map:
\[
\text{ev}_p: k[x_1, \dots, x_n] \to k \qquad f \mapsto f(p)
\]
is an algebra homomorphism whose kernel, $m_p = ( x_1-a_1, \dots, x_n-a_n )$, is a maximal ideal. The quotient $k[x_1, \dots, x_n]/m_p \cong k$ recovers the value of the function at that point.

We have established a similar setup for a finite abelian group $G$, with the convolution algebra
$(\Hom_\Top(G, \mathbb{C}), *, +)$ taking the role of the function algebra. A natural question arises: can we
recover the ``points'' of our space from the maximal ideals of this algebra? Let's first attempt a direct
analogy. We might assume the points are the elements of the group $G$ itself. For a point $g_i \in G$, we
can define the point-evaluation map, which is a linear map:
\[
\text{ev}_{g_i}: \Hom_\Top(G, \mathbb{C}) \to \mathbb{C} \qquad f \mapsto f(g_i)
\]
The kernel of this map is the vector subspace of functions that are zero at $g_i$. However, for this subspace to be an ideal, it must be closed under convolution. Let $f \in \ker(\text{ev}_{g_i})$ (so $f(g_i)=0$) and let $h$ be any function in our algebra. Consider their convolution:
\[
(h * f)(g_i) = \sum_{k \in G} h(k)f(k^{-1}g_i)
\]
Just because $f(g_i)=0$ does not imply that the sum on the right-hand side is zero. The value of $f$ at other points, $f(k^{-1}g_i)$, can be non-zero. Therefore, the kernel of the point-evaluation map is \emph{not} an ideal with respect to convolution. This implies the points of $G$ do not correspond to the maximal ideals of its convolution algebra.

The key insight is that the correct ``geometric points'' for the convolution algebra are not the elements of $G$, but the characters of the dual group $\widehat{G}$.

\begin{prop}{Characters and Maximal Ideals}{charsAndMaxIdeals}
    Let $G$ be a finite abelian group. The maximal ideals of the convolution algebra $(\Hom_\Top(G,
	\mathbb{C}), +, *)$ are in bijective correspondence with the characters $\chi \in \widehat{G}$.
\end{prop}

\begin{Proof}
    For each character $\chi \in \widehat{G}$,  define a Fourier evaluation map:
    \[
        \widehat{\text{ev}}_\chi: \Hom_\Top(G, \mathbb{C}) \to \mathbb{C} \qquad f \mapsto \hat{f}(\chi)
    \]
    This map takes a function and returns its Fourier coefficient at the frequency $\chi$.
	By~\ref{th:fourierReciprocity} $\widehat{\text{ev}}_\chi$ is a $\C$-algebra homomorphism. Since
	$\widehat{\text{ev}}_\chi$ is a surjective algebra homomorphism onto a field ($\mathbb{C}$), its kernel,
	which we denote by $I_\chi$, is a maximal ideal:
    \[
        I_\chi = \ker(\widehat{\text{ev}}_\chi) = \{ f \in \Hom_\Top(G, \mathbb{C}) \mid \hat{f}(\chi) = 0 \}
    \]
    By the First Isomorphism Theorem for rings, we have:
    \[
        \Hom_\Top(G, \mathbb{C}) / I_\chi \cong \mathbb{C}
    \]
	This shows that every character $\chi \in \widehat{G}$ defines a maximal ideal. Conversely, After applying
  the Fourier transform we get the usual multiplication and ring structure, and then an approach similar to
  the Gelfand-Kolmogorov theorem gives the desired result.


  For the infinite case, the maximal ideals of the Banach algebra $L^1(G, \mu)$ arise in this same way from
  homomorphisms into $\mathbb{C}$ (see the \emph{Gelfand-Mazur Theorem}).
\end{Proof}

Thus, with the convolution algebra $\mathbb{C}[G]$ as our primary object, the natural geometric space that emerges
is not the group $G$ itself, but its Pontryagin dual, $\widehat{G}$. The Fourier transform is the mechanism
that allows us to evaluate the functions of our algebra at the points of this ``true'' underlying space. This
perspective, generalized by Gelfand, is the foundation of non-commutative geometry.


\begin{center}
\begin{tabular}{l|l}
	\textbf{Concepts} & \textbf{Classical Algebraic Geometry} \\
	\hline
	Algebra & Polynomial Ring $(k[x_1, \dots, x_n], \cdot)$ \\
	Product & Polynomial Multiplication \\
	Geometric Space & Affine Space $k^n$ \\
	A "Point" $p$ & A tuple $(a_1, \dots, a_n) \in k^n$ \\
	Evaluation Map & $\text{ev}_p: f \mapsto f(p)$ \\
    Maximal Ideal for $p$ & $\ker(\text{ev}_p) = \langle x_i - a_i \rangle$ \\
\end{tabular}

\vspace{1em}

\begin{tabular}{l|l}
	\textbf{Concepts} & \textbf{Harmonic Analysis on G} \\
	\hline
	Algebra & Group Algebra $(\mathbb{C}[G], *)$ \\
	Product & Convolution \\
	Geometric Space & The Dual Group $\widehat{G}$ \\
	A "Point" $p$ & A character $\chi \in \widehat{G}$ \\
	Evaluation Map & Fourier ``Evaluation'' $\widehat{\text{ev}}_\chi: f \mapsto \hat{f}(\chi)$ \\
    Maximal Ideal for $p$ & $\ker(\widehat{\text{ev}}_\chi) = I_\chi$ \\
\end{tabular}
\end{center}




% \newpage
\printbibliography

\end{document}
